\chapter{Specifying events of interest}\label{ch:specify}


\section{Focus event of this project}

In this thesis we explicitly set the focus on the heatwave event (TODO: state this in the intro directly).
However, the method we develop is variable agnostic, and could be used and tuned to generate other extreme events of interest following a very similar logic.

We narrow the focus on purpose with the purpose of keeping the need for expert evaluation as far as possible from this project, 
in which we want to focus on the methodology itself.

\section{Mask}

In our setup we are interested in generating a heatwave scenario in a specific region of interest.
This means, we want to guide the base forecast to generate higher temperatures over the selected region.

We define the region of interest as a simple bounding box, which defines a set of weigths (sum to 1) over a tensor having the same dimension as the weather states, non-zero at the variable and region of interest.

It is really important to understand how to evaluate what we are doing here.
Since we are aiming at a specific region of interest, and we are guiding only wrt the region of interest, 
we are loosing the global coherence of the weather states.
We shall thus evaluate the ability of our method to generate locally coherent weather scenarios, over the region of interest.

\section{Target trajectory}

In the following we define a climatological heuristic to define the target trajectory.

In order to define plausible average trajectories of change, we define following procedure:
1. Select start state
2. Define mask
3. Retrieve historical states that are consistent with the start state and the mask (+- some tolerance)
4. Compute the distribution of average trajectories of change for the historical states
5. Filter a quantile of the distribution to define the target trajectory (e.g. 0.95 quantile for extreme events)
6. Define a gaussian distribution at each time step of the target trajectory, with mean equal to the target trajectory and variance equal to the variance of the historical states at that time step.
7. Sample enough trajectories from the gaussian distribution to define a set of target trajectories that will be used to guide the generative model.
8. Select the target trajectories of interest in terms of extremeness to build scenarios.
9. Save the trajectories and corresponding states used to define the sampling distribution. These will be used to evaluate the plausibility of the generated trajectories as means of comparison.
10. Run benchmark tests and evaluate.

This heuristic enables us to anchor the guidance target and define a benchmark at the same time.
More concretly, we avoid the issue of having to tune the strength of the guidance to achieve a certain plausible result.
We reverse the problem by defining the plausible result first and then designing an algorithm to achieve that result and checking that both the algorithm and its result are well behaved.

We explicitly define a heatwave in order to design guidance methods. 
This means that we decide to represent a specific subset of state trajectoriesin this work.
However, our method is flexible to any kind of target: you can swap your definition of the event you want to replicate and use the guidance to generate that kind of event.
This really only affects the target definition.

This choice enables us to side step the problem that guiding towards a classifier's gradients or minimizing / maximizing a certain variable and stopping when we achieved the wanted result.
If we specify the carachteristic of the wanted result we side step this problem entirely.
The problem then shifts to retrieve a representative set of the event we have in mind.
This is however a requirement for the first methodology.
So we could say that the problem shifts towards defining an appropriate algorithm that can achieve a representative state similar to the specified set.
Making this step explicit directly enables us to then compare the results against this set,
and characterize its flaws / benefits directly.

Note that this whole pipeline enables us to sample trajectories wrt to the historical data.
However, we can define arbitrary targets for our method, which we will do in chapter 6 when testing the method.
In chapter 6.3 we will explore the usage of this heuristic detached from the method evaluation, in the context of the use cases we are actually interested in generating.

\section{Pipeline summary}

In summary we are explicitly defining our event of interest in terms of the average temperature over a specified area of interest, w.r.t. the climatology over the same area from the year 2020 to today.

--- the concretization of our approach maybe go under experiments?

--- maybe we don't need a summary section in the first place.

--- put here a viz